%% kapa: exam=1 round=28 year=2017 period=1 form=A subject=civil_law subject_ko=민법 seq=18 number=18
%% source: Q-Net 감정평가사 기출문제 (cst00309) artlSeq=5199603 file=extracted/1008/5199603/2017년 제28회 감정평가사 1차시험 문제지/1교시 A형.pdf
%% conversion: kordoc 4.24.51 -> tools/gen1.py (auto); review=proofread-sonnet; flags=none
\begin{question}{18}
조건에 관한 설명으로 옳지 않은 것은? (다툼이 있으면 판례에 따름)
\begin{choices}
  \item 조건은 법률행위의 효력의 발생 또는 소멸을 장래의 불확실한 사실의 성부에 의존하게 하는 법률행위의 부관이다.
  \item 불능조건이 해제조건이면 조건 없는 법률행위가 된다.
  \item 조건의사가 있더라도 법률행위의 내용으로 외부에 표시되지 않은 경우, 그것만으로는 법률행위의 조건이 되지 않는다.
  \item 부관이 붙은 법률행위에 있어서 부관에 표시된 사실의 발생 유무에 상관없이 그 채무를 이행해야 하는 경우에는 조건으로 보아야 한다.
  \item 정지조건부 법률행위의 경우에는 조건성취로 권리를 취득하는 자가 조건성취 사실에 대한 증명책임을 진다.
\end{choices}
\end{question}
\answer{4}
